\clearpage\begingroup\NoSourceHead\SourceLeaf{220}
\makeatletter\tagsleft@true\let\veqno\@@leqno\makeatother
\begin{center}
[III.]\par\medskip
\textbf{BERECHNUNG DES NEUMONDS TISRI FÜR JEDES JÜDISCHE JAHR $A$.}\par\medskip
\rule{\textwidth}{.35pt}\par
{\small [Handschriftliche Eintragung in \textsc{Christian Wolf}, Elementa matheseos universae, tomus IV. --- Von \textsc{Gauss} 1800 erworben.]}\par
\rule{\textwidth}{.35pt}
\end{center}
\begin{center}[1.]\end{center}
\noindent\hspace*{\parindent}Den Tag von Molad Tohu oder den 7. Oct. des Jahres 953 der Julian. Periode bezeichnen wir mit 2, den folgenden mit 3 etc. in inf. Man fragt, auf welchen Tag $x$ der Neumond Tisri des Jahres $A$ fallen werde.

Auflös[ung]. Es entstehe aus der Division der Zahl $7A+13$ mit $19$ der Rest $a$, so ist
\[
\begin{aligned}
x&=\frac{5604}{25\,920}+\frac{1}{19}(235A-234-a)\times\left(29\frac{13753}{25\,920}\right)\\
&=A\times365\frac{24311}{98\,496}-1\frac{272\,953}{492\,480}a-363\frac{234\,606}{492\,480}.
\end{aligned}
\]
\begin{center}[2.]\end{center}
\noindent\hspace*{\parindent}Auf welchen Tag des Julianischen Kalenders fällt der Neumond Tisri des Jahres $A$?

Jahr Christi $=A-3761$.

$A-1$ mit $4$ dividirt, gebe den Rest $b$, so ist der gesuchte Tag
\[
A-3761\ \text{Jahr,}\qquad \frac{1}{4}b-\frac{313}{98\,496}A-1\frac{272\,953}{492\,480}a+8\frac{380\,994}{492\,480}\ \text{Octob.}
\]

\clearpage\SourceLeaf{221}\HistoricalHead{216}{L}{CHRONOLOGIE.\quad NACHLASS.}
\begin{center}\small BEMERKUNGEN.\end{center}
\begingroup\small
\noindent\hspace*{\parindent}Dieser Ansatz unterscheidet sich in zweifacher Weise von \textsc{Gauss}' \textit{Bestimmung des jüdischen Osterfestes} (Werke VI, S. 80, vergl. auch X\textsubscript{1}, S. 560). Erstens wird statt des bei der Osterberechnung auftretenden Restes von $12A+17$ durch $19$ der für die Rechnung bequemere von $7A+13$ durch $19$, wofür man den gleichwertigen von $7A-6$ durch $19$ setzen kann, verwendet. Zweitens bildet, was bedeutsamer ist, statt des jüdischen Osterfestes --- das \textsc{Gauss} in seiner erwähnten Veröffentlichung wohl nur in Analogie zu seiner früher gegebenen Formel für das christliche Osterfest behandelt hat --- hier der Neumond Tišri\begingroup\renewcommand{\thefootnote}{*)}\footnote{Tišri ist der siebente Monat des jüdischen Jahres; gemäß Leviticus Kap. 23, Vers 24 wird (abgesehen von gewissen genau festgelegten Ausnahmen) an seinem Neumond der »Feiertag des Blasens«, das jüdische Neujahrsfest, begangen, mit dessen Eintritt die Zählung des neuen Jahres beginnt.}\endgroup, der das wirkliche Fundament der jüdischen Chronologie ist, den Ausgangspunkt. Es soll hier eine Bestätigung der beiden von \textsc{Gauss} aufgestellten Ausdrücke gegeben werden; im Jahresbericht der Deutschen Mathematiker-Vereinigung 26, 1917, S. 304 zeige ich, wie man von dem \textsc{Gauss}schen Ansatze aus zu einer formelmäßigen Bestimmung des jüdischen Neujahrsfestes und des jüdischen Osterfestes gelangen kann.

Der jüdische Kalender hat in einem Zyklus von $19$ Jahren immer je $7$ Schaltjahre von $13$ Monaten, nämlich die Jahre $3,6,8,11,14,17$ und $19$ des Zyklus, während die andern Jahre aus $12$ Monaten bestehen. Geht man von dem Beginn des 9. Jahres irgend eines Zyklus, also von dem Beginn des Jahres $19Z+9$, wobei $Z$ irgend eine ganze positive Zahl bedeutet, immer um $8$ volle Jahre zurück, so sind, wie man sich überzeugt, in jedem solchen achtjährigen Intervall stets $3$ Schaltjahre enthalten. Nun läßt sich jede ganze positive Zahl $A$ in einer und auch nur in einer einzigen Weise in die Form
\begin{equation}
A=19Z+9-8a\tag*{(1)}
\end{equation}
\noindent bringen, wobei $a$ und $Z$ ganzzahlig sind und $a$ der Ungleichung $0\leqq a<19$ genügt. Setzt man die mit $7$ multiplizierte Gleichung (1) in die Form
\begin{equation}
7A=19(7Z-3a+3)+a+6,\tag*{(1*)}
\end{equation}
\noindent so findet man sofort ihre Lösungen
\begin{equation}
a=R\left(\frac{7A-6}{19}\right)\tag*{(2)}
\end{equation}
\noindent und
\begin{equation}
Z=\frac{Q\left(\frac{7A-6}{19}\right)+3R\left(\frac{7A-6}{19}\right)-3}{7},\tag*{(3)}
\end{equation}
\noindent wo $Q,R$ den ganzzahligen nicht negativen Quotienten bezw. Rest bedeuten. Aus der Darstellung der ganzen positiven Zahl $A$ in der Form (1) ergibt sich unmittelbar die Anzahl der Schaltjahre, die von Beginn der jüdischen Zeitrechnung, dem Beginn des jüdischen Jahres 1, bis zum Anfang des jüdischen Jahres $A$ vergangen ist. Von Beginn der jüdischen Zeitrechnung bis zu Beginn des Jahres $19Z+9$ sind $7Z+3$ Schaltjahre verflossen, nämlich $7Z$ Schaltjahre in den völlig abgelaufenen $Z$ Zyklen und die Schaltjahre 3, 6, 8 der acht verflossenen Jahre des $(Z+1)^{\mathrm{ten}}$ Zyklus. Geht man vom Beginn des Jahres $19Z+9$ bis zum Anfang des Jahres $19Z+9-8a$ zurück, so liegen in diesen $a$ achtjährigen Intervallen, die wir zurückgingen, nach der oben gemachten Bemerkung $3a$ Schaltjahre. Mithin beträgt die seit Beginn der jüdischen Zeitrechnung bis zum Anfang des jüdischen Jahres $A=19Z+9-8a$ verflossene Anzahl von Schaltjahren $7Z+3-3a$; diese Zahl ist nach (3) und (2) gleich $Q\left(\frac{7A-6}{19}\right)$. Die bis zum Beginn des jüdischen Jahres
\clearpage\SourceLeaf{222}\HistoricalHead{217}{R}{BEMERKUNGEN.}\HistoricalFoot{XI\textsubscript{1}.}{28}
\noindent $A$ vergangene Anzahl von Schaltjahren ist daher gleich der größten in $\dfrac{7A-6}{19}$ enthaltenen ganzen Zahl $Q\left(\dfrac{7A-6}{19}\right)$.

Die bis zu Beginn des Jahres $A$ vergangenen $A-1$ Jahre enthalten $12(A-1)+Q\left(\frac{7A-6}{19}\right)$ Monate, nämlich erstens $12(A-1)$ reguläre Monate und zweitens $Q\left(\frac{7A-6}{19}\right)$ Schaltmonate.

Nach den Gleichungen (3) und (2) ist
\begin{equation}
12(A-1)+Q\left(\frac{7A-6}{19}\right)=12(A-1)+7Z+3-3a.\tag*{(4)}
\end{equation}
\noindent Nun ist nach der Gleichung (1*)
\[
7Z+3-3a=\frac{7A-a-6}{19};
\]
\noindent folglich geht der Ausdruck (4) über in
\[
\begin{aligned}
12(A-1)+\frac{7A-a-6}{19}
&=\frac{1}{19}[12\mathbin{.}19(A-1)+7A-a-6]\\
&=\frac{1}{19}(235A-234-a).
\end{aligned}
\]

Nach jüdischer Annahme fiel der Neumond Tišri des Jahres 1, der Molad Tohu\begingroup\renewcommand{\thefootnote}{*)}\footnote{Im Hebräischen bedeutet \textls[120]{Molad} soviel wie Geburt, hier insbesondere das Erscheinen des Mondlichts, also \textls[120]{Neumond}, und \textls[120]{Tohu} soviel wie \textls[120]{wüst} mit Bezug auf die Schöpfung, vergl. Genesis Kap. 1, Vers 2. Molad Tohu ist der Neumond der Schöpfung, der Ausgangspunkt der jüdischen Zeitrechnung.}\endgroup, auf einen Montag und zwar traf er $5\frac{204}{1080}$ Stunden $=\frac{5604}{25920}$ Tage nach Beginn des fraglichen Montags ein; als Tagesanfang gilt hierbei nach jüdischer Zählung die Zeit 6 Uhr abends, also eine um 6 Stunden frühere Zeit als nach christlicher Rechnung. Da der jüdische Monat $29\frac{13753}{25920}$ Tage besitzt, sind seit Beginn jenes Tages, an dem der Molad Tohu stattfand, wie \textsc{Gauss} angibt,
\begin{equation}
\begin{aligned}
&\frac{5604}{25920}+\frac{1}{19}(235A-234-a)\mathbin{.}\left(29\frac{13753}{25920}\right)\\
&\qquad=365\frac{24311}{98496}A-1\frac{272953}{492480}a-363\frac{234606}{492480}
\end{aligned}\tag*{(5)}
\end{equation}
\noindent Tage verflossen. Dabei gilt als Tagesbeginn des Molad Tohu, der auf den 7. Oktober 3761 v. Chr.\begingroup\renewcommand{\thefootnote}{**)}\footnote{Das Jahr 1 der sogenannten julianischen Periode ist das 4713. v. Chr., so daß die von \textsc{Gauss} oben S. 215 gegebene Bestimmung »des Jahres 953 der Julian. Periode« auf das Jahr $4714-953=3761$ v. Chr. hinaus kommt.}\endgroup\ fiel, nach jüdischer Rechnung der voraufgehende Sonntag abends 6 Uhr.

\textls[120]{Man setze den Ausdruck (5) gleich $x+\xi$, wobei $x$ den ganzzahligen Bestandteil bezeichnet und $\xi$ der positive echte Bruch ist.} Numeriert man alsdann Samstag den 5. Oktober 3761 v. Chr. mit 0, den folgenden Sonntag mit 1, Montag, den Tag von Molad Tohu, mit 2, den folgenden Dienstag mit 3 u.s.w., so erhält der Neumond Tišri des Jahres $A$ die Nummer $x+2$. Die \textsc{Gauss}sche Angabe über die Numerierung ist also dementsprechend abzuändern.

Um das Datum zu bestimmen, auf das der Neumond Tišri des jüdischen Jahres $A$ fällt, verfahren wir folgendermaßen: Wie wir gesehen haben, fällt dieser Neumond auf den $(x+2)^{\mathrm{ten}}$ Tag nach dem 5. Oktober 3761 v. Chr., d. h. auf den $(x+7)^{\mathrm{ten}}$ Tag nach dem 0. Oktober 3761 v. Chr.; er möge nun $A-1$ julianische Jahre und $t$ Tage nach dem 0. Oktober 3761 v. Chr. fallen, $t$ bedeute hierbei im Falle, daß der Neumond Tišri vor dem 0. Oktober eintrifft, eine negative Zahl. Die $A-1$ julianischen Jahre besitzen $Q\left(\frac{A-1}{4}\right)$ Schalttage; denn das julianische Jahr 3761 v. Chr., von dessen 0. Oktober gezählt werden soll, ist als
\clearpage\SourceLeaf{223}\HistoricalHead{218}{L}{CHRONOLOGIE.\quad NACHLASS.}
\noindent Schaltjahr anzusehen, dessen Schalttag bereits vorüber ist, und auf dieses Jahr folgen die Schaltjahre immer in vierjährigen Intervallen, so daß das erste derselben nach Ablauf von $A-1=4$ julianischen Jahren zu berücksichtigen ist. Die fraglichen $A-1$ julianischen Jahre enthalten daher
\[
365(A-1)+Q\left(\frac{A-1}{4}\right)=365\frac{1}{4}(A-1)-\frac{1}{4}R\left(\frac{A-1}{4}\right)
\]
\noindent Tage; denn es ist
\[
A-1=4Q\left(\frac{A-1}{4}\right)+R\left(\frac{A-1}{4}\right).
\]
\noindent Mithin fällt der Neumond Tišri des Jahres $A$ genau
\[
t+365\frac{1}{4}(A-1)-\frac{1}{4}R\left(\frac{A-1}{4}\right)
\]
\noindent Tage nach dem 0. Oktober 3761 v. Chr. Da er aber, wie oben gezeigt, $x+7$ Tage nach dem genannten Termin fällt, so hat man die Gleichung
\[
t+365\frac{1}{4}(A-1)-\frac{1}{4}R\left(\frac{A-1}{4}\right)=x+7
\]
\noindent oder
\begin{equation}
t=x+7-365\frac{1}{4}(A-1)+\frac{1}{4}R\left(\frac{A-1}{4}\right).\tag*{(6)}
\end{equation}
\noindent Addiert man bei (6) rechts und links $\xi$ und ersetzt rechter Hand $x+\xi$ durch seinen Wert nach (5), so wird
\begin{equation}
\begin{aligned}
t+\xi={}&365\frac{24311}{98496}A-1\frac{272953}{492480}a-363\frac{234606}{492480}\\
&+7-365\frac{1}{4}(A-1)+\frac{1}{4}R\left(\frac{A-1}{4}\right)
\end{aligned}\tag*{(7)}
\end{equation}
\noindent oder
\begin{equation}
t+\xi=\frac{1}{4}R\left(\frac{A-1}{4}\right)-\frac{313}{98496}A-1\frac{272953}{492480}a+8\frac{380994}{492480}\text{\textsuperscript{*)}}.\tag*{(7*)}
\end{equation}
\begingroup\renewcommand{\thefootnote}{}\footnotetext{*) Erweitert man rechts vom Gleichheitszeichen Zähler und Nenner des Koeffizienten von $A$ mit 5, so sieht man, daß $t+\xi$ auf den Nenner $492480=19\mathbin{.}24\mathbin{.}1080$ gebracht werden kann. Es ist aber für die Ausführung und Nachprüfung der Rechnung von Bedeutung, daß sich $\xi$ stets als echter Bruch mit dem Nenner $25920=24\mathbin{.}1080$ schreiben läßt. $\xi$ kann nämlich auch als der echte Bruch des in (5) links vom Gleichheitszeichen stehenden Ausdrucks erklärt werden, bei dem $\dfrac{235A-234-a}{19}$ als Anzahl der bis zum Beginn des Jahres $A$ verflossenen Monate eine ganze Zahl ist.}\endgroup
\noindent Es ist mithin $t$ der ganzzahlige Bestandteil, $\xi$ der positive echte Bruch des Ausdrucks, der auf der rechten Seite von (7*) steht. Dieses Resultat, wonach also der Neumond Tišri des jüdischen Jahres $A$ auf den $t^{\mathrm{ten}}$ julianischen Oktober des Jahres $A-3761$ n. Chr. fällt, stimmt mit der Angabe von \textsc{Gauss} überein. Ist $t$ gleich Null oder negativ so ist der $t^{\mathrm{te}}$ julianische Oktober mit dem $(30+t)^{\mathrm{ten}}$ julianischen September gleichbedeutend.
\begin{flushright}\textsc{Loewy.}\end{flushright}
\endgroup
\begin{center}\rule{.25\textwidth}{.35pt}\end{center}
\clearpage\endgroup
